\begin{figure}[h]
\begin{minipage}{\textwidth}
  \includegraphics[width=\linewidth]{bindle/bindle-cupy-traits.py/bindle/teeplots/cupy-traits/errorbar=ci+exclude=1D-demes+hue=population-structure+kind=line+palette=dark2+row=population-size+style=initial-conditions+viz=relplot+x=available-beneficial-mutations+y=fix-prob+ext=.pdf}%
\end{minipage}

\begin{minipage}{\textwidth}
  \caption{%
    \textbf{Population structure increases nonmutator resilience to available adaptive potential.}
    \footnotesize
    Lineplot strips show relationship between adaptive potential and mutator fixation probability across surveyed population sizes, with error bands indicating bootstrapped 95\% confidence intervals.
    In the 50/50 treatment, experiments were initialized with an even mix of mutator and nonmutator agents.
    In the \textit{de novo} treatment, mutators are initially absent and arise spontaneously \textit{de novo} with probability $10^{-6}$.
    Effects of population structure are more pronounced under \textit{de novo} conditions, when initial supply of mutators is scarce.
    Under well-mixed conditions, increase in population size favors mutators; nonmutators reliably persist in very large populations through fewer than 10 available beneficial mutations.
    In contrast, under 2D deme structure, very large population sizes are favorable to nonmutators.
    With such strong spatial structure, very large populations of nonmutators resist mutator fixation even when upwards of 10 beneficial mutations are available.
    Results are under GPU simulation, using counter-based genome model.
    Supplementary \cref{fig:fixheat-5050-cupy,fig:fixheat-denovo-cupy} detail results in a tabular format.
    [Dashed black lines show the analytic approximation of Section \ref{sec:model}.]
  }
  \label{fig:spatial-structure-combined}
\end{minipage}
\end{figure}
